\subsection{Why Hierarchical Structure Provides Stability}

The stability advantage of recursive bisection derives from its hierarchical decomposition. We propose the following mechanism:

\textbf{Fundamental Geographic Splits:} Top-level bisections naturally align with major geographic features—mountain ranges, rivers, urban-rural divides, historical regions (e.g., Alabama's Black Belt). These fundamental divisions persist across decades because they reflect geological and historical constants.

\textbf{Localized Adjustments:} When demographics shift within a region, recursive bisection adjusts lower-level splits without affecting the top-level structure. For example, if District 1's population grows, adjustments occur between Districts 1 and 2 (both in Region A) but don't propagate to Region B.

\textbf{Cascade Prevention:} N-way's flat structure lacks this containment property. Population growth in District 1 can trigger boundary changes that cascade through Districts 2, 3, \ldots, $k$ as the algorithm re-optimizes all districts simultaneously to maintain balance.

\textbf{Alabama Case Study:} Alabama's [3,4] top-level split likely divides the Black Belt (predominantly minority counties in west-central Alabama) from the rest of the state. This geographic-demographic feature persisted from 2010 to 2020, enabling recursive bisection to preserve the split. Within each region, lower-level adjustments accommodated population changes without disturbing the fundamental division.

\subsection{The Performance-Stability Tradeoff}

Our findings, combined with prior work on VRA performance (Paper 2), reveal a fundamental tradeoff:

\begin{center}
\begin{tabular}{lcc}
\toprule
Criterion & N-Way Advantage & Recursive Advantage \\
\midrule
2020 Minority Concentration & \textbf{+4.3 pts} & — \\
2020-2030 Stability & — & \textbf{+10-13 pts} \\
Runtime & \textbf{-67\% faster} & — \\
Interpretability & — & \textbf{Tree structure} \\
\bottomrule
\end{tabular}
\end{center}

\textbf{Decision Framework:} Jurisdictions should choose methods based on context:

\begin{enumerate}
    \item \textbf{Borderline VRA Compliance:} If achieving one additional minority-majority district is marginal (e.g., Alabama's potential second district), use n-way for maximum immediate performance.

    \item \textbf{Comfortable VRA Margins:} If multiple minority-majority districts are reliably achievable, use recursive bisection to minimize 2030 disruption. The +10-13 point stability gain outweighs the -4.3 point performance cost.

    \item \textbf{Rapid Prototyping:} For exploratory analysis or large-scale sweeps, n-way's speed advantage (67\% faster) dominates.

    \item \textbf{Public Transparency:} Recursive bisection's interpretable tree structure may aid public understanding and legal defensibility.
\end{enumerate}

\textbf{Hybrid Approach:} One could initialize 2030 redistricting with 2020's recursive tree structure, adjusting only lower levels. This would preserve top-level splits while accommodating demographic changes—potentially combining stability and performance.

\subsection{Implications for 2030 Redistricting}

Our results generate testable predictions for the 2030 census cycle:

\textbf{Prediction 1:} States that used recursive bisection in 2020 will experience less disruption in 2030 redistricting (measured by tract reassignment rates from 2020 to 2030).

\textbf{Prediction 2:} Preserving 2020's top-level split in 2030 will yield higher stability than re-running recursive bisection from scratch.

\textbf{Prediction 3:} The stability advantage will correlate with demographic volatility—states with rapid demographic change (e.g., Georgia, North Carolina) will benefit more from recursive methods than stable states (e.g., Vermont).

When 2030 census data arrives, researchers can test these predictions directly, validating or refuting our mechanistic explanation.

\subsection{Generalization Beyond Redistricting}

The performance-stability tradeoff extends to any domain requiring periodic re-partitioning:

\textbf{School Districts:} As enrollment shifts, school boundaries must adjust every 5-10 years. Hierarchical methods could minimize student reassignment.

\textbf{Service Territories:} Hospitals, fire departments, and police precincts rebalance catchments as populations grow. Stability reduces confusion and maintains institutional knowledge.

\textbf{Sales Regions:} Companies periodically rebalance sales territories. Hierarchical methods preserve regional organization while adjusting local boundaries.

\textbf{General Principle:} When re-partitioning frequency is low relative to inter-partition timescale, stability becomes valuable. Hierarchical methods provide continuity through structural preservation, while flat methods optimize each instance independently at the cost of disruption.

\subsection{Limitations of Current Analysis}

\textbf{Two Cycles Only:} We analyze 2010-2020 but lack 2000 data for three-cycle validation. Adding 2000 would strengthen claims about long-term stability.

\textbf{Algorithm-Specific:} Results apply to METIS-based partitioning. Other algorithms (spectral clustering, simulated annealing) may exhibit different stability properties.

\textbf{Tract Boundary Changes:} Census Bureau relationship files imperfectly map 2010-2020 tract changes. Some splits/merges introduce noise in stability metrics.

\textbf{Five States:} Limited to southern states with VRA scrutiny. Generalization to all 50 states requires additional experiments.
